\section{الفصل~\ref{sec:eetypes}. العصبونات أجهزةً}

أوضاع عمل الخلايا المفردة مُقاسة مباشرة، بالتيار عبر القطب وتسجيل الجهد؛ ورياضيات المكامل والتقسيم إلى صنفين نظرية كلاسيكية؛ أما استعمال الدماغ لإعادة تهيئة الأوضاع في الحساب فيبقى فرضية.

{\footnotesize
\setlength{\tabcolsep}{3pt}\renewcommand{\arraystretch}{1.15}
\begin{longtable}{>{\raggedright}p{0.5cm}>{\raggedright}p{4.0cm}>{\raggedright}p{5.6cm}>{\raggedright}p{2.3cm}>{\raggedright\arraybackslash}p{1.7cm}}
\toprule
م & القضية & التجربة وما قيس & المصدر & الحالة \\
\midrule\endfirsthead
\toprule
م & القضية & التجربة وما قيس & المصدر & الحالة \\
\midrule\endhead
1 & الرنّان: تيار بطيء يعاكس تغير الجهد يجعل العصبون مرشح تمرير نطاقي قمته عند بضعة هرتزات إلى عشراتها (القسم~\ref{sec:resonator}) & في الشرائح حُقن في عصبونات \nt{GLUT} تيار جيبي يزداد تردده تدريجيًا وقيست الممانعة $|Z(f)|$: القمة عند $2\text{--}5\,\text{Hz}$ في $33^\circ$C (ونحو $7\,\text{Hz}$ في $38^\circ$C)؛ ويصنع الرنين تياران بطيئان، بوتاسيومي وكاتيوني، ويقويه تيار صوديوم مستمر. وتعرض المراجعة الحيلة نفسها لأصناف كثيرة من الخلايا & \cite{HuVervaekeStorm2002,HutcheonYarom2000} & مُقاس \\
2 & التيار البطيء يتصرف كالمحاثة، ويشكّل مع سعة الغشاء دارة رنين & قيست الاستجابة دون العتبية لمحوار على التيار ووُصفت بمعادلات الموصليات بعد تقريبها خطيًا؛ والدارة المكافئة تحتوي محاثة ”ظاهراتية“ تعتمد قيمتها على الجهد & \cite{Mauro1970} & نظرية \\
3 & الثابت الزمني لمجموعة بتغذية راجعة $\tau_{\text{eff}}=\tau/(1-w)$~\eqref{eq:perfectint}؛ عند $\tau=0.1$ ثانية و$w=0.995$ يكون $20$ ثانية & مشتق في الفصل من المعادلة الخطية للمجموعة؛ واقتُرح آليةً لحفظ موضع العين في عمل نظري & اشتقاق في الفصل؛ \cite{Seung1996} & نظرية \\
4 & مكامل موضع العين يحفظ المدخل بالشبكة لا بخاصية خلية مفردة & تسجيل داخل خلوي عند سمكة في العقدة التي تحفظ موضع العين: بعد دوران سريع يتغير تردد العصبون درجةً ويبقى؛ والتيار القصير في خلية واحدة لا يسبب إلا إزاحة عابرة، أما درجات الجهد فتبقى حتى دون العتبة، إذ تحفظها المداخل المشبكية & \cite{Aksay2001} & مُقاس \\
5 & المكامل غير التسريبي يحتاج إلى ضبط دائم بالتعلم؛ وعند اختلال الضبط إما ينسى وإما ينجرف إلى الإشباع & دُرّبت سمكة بتغذية راجعة بصرية تتناسب مع $\pm$ موضع العين: صار الحفظ غير مستقر أو تسريبيًا، وهبط الثابت الزمني للعصبونات بأكثر من رتبة، إلى $<1$ ثانية؛ وأعادت التغذية الراجعة البصرية العادية الاستقرار تدريجيًا & \cite{Major2004} & مُقاس \\
6 & العصبون ثنائي الاستقرار: مدخل قصير يشغّل هضبة طويلة، والتثبيط يطفئها؛ والسيروتونين يشغّل هذه الخاصية (القسم~\ref{sec:bistable}) & تسجيل داخل خلوي لعصبونات \nt{ACET} التي تتحكم في العضلات عند القط: إثارة مشبكية قصيرة تنقل العصبون إلى عمل طويل، وتثبيط قصير يعيد ضبطه؛ وعند القط مقطوع الحبل الشوكي تظهر الحالة بعد إعطاء سلائف السيروتونين & \cite{Hounsgaard1988} & مُقاس \\
7 & صنفان: المكاملات (يرتفع التردد من الصفر) والرنّانات (عند العتبة تردد منتهٍ فورًا) & الصنفان مشتقان من نظرية التشعبات. والتجربة: في شرائح القشرة تبدأ عصبونات \nt{GLUT} العمل بتردد صغير كيفما شئنا (النوع~1)، وعصبونات \nt{GABA} السريعة تبدأ بقفزة من تردد منتهٍ (النوع~2) & \cite{Izhikevich2007,Tateno2004} & مُقاس \\
8 & العصبون الواحد مكامل أو كاشف تزامن: التثبيط يقصّر نافذة الجمع & في شرائح الحُصين النافذة التي تُجمع فيها المداخل المثيرة حتى العتبة أقصر من $2\,\text{ms}$؛ ويقصّرها التثبيط الأمامي الذي يصل فور الإثارة؛ وفي التغصنات، حيث التثبيط أضعف، النافذة أوسع & \cite{Pouille2001} & مُقاس \\
9 & خط تأخير من المحاوير (الجدول~\ref{tab:eetypes}) & قيست عند البومة الصيادة التأخيرات في المحاوير الذاهبة إلى كواشف التزامن: اختلاف طول المحوار يعطي تأخيرًا مختلفًا، والكواشف مضبوطة على فرق زمن وصول الصوت & \cite{Carr1990} & مُقاس \\
10 & ناظمة إيقاع في اليقظة وخلية صامتة في النوم & عصبونات \nt{SERO} تعمل بانتظام تقريبًا نهارًا، وتتباطأ في النوم البطيء، وتكاد تصمت في النوم السريع (التفاصيل في الفصل~\ref{sec:sleep}) & \cite{JacobsAzmitia1992,McGintyHarper1976} & مُقاس \\
11 & يحدد الجهازَ القنواتُ الأيونية وقنواتُ البث التي تضبطها (القضية~\ref{prop:eetypes}) & ثبت على شبكات صغيرة: المواد المنظِّمة تغيّر الخصائص الذاتية للعصبونات وقوة المشابك، فتُصدر الدارة نفسها إيقاعات مختلفة & \cite{Marder2012} & مُقاس \\
12 & الدماغ يستعمل إعادة تهيئة الأوضاع في الحساب (القضية~\ref{prop:eetypes}) & لا توجد تجربة مباشرة كانت فيها إعادة تهيئة وضع الخلية لازمة لحل مسألة؛ ثمة فقط تجارب تغيّر فيها المنظِّمات مخرج الدارة & \cite{Marder2012} & فرضية \\
\bottomrule
\end{longtable}}
